% ============================================================
%  Section 1, Segment 1A — Finger Tapping / 点弦
%  原书页码 4-13
% ============================================================
\chapter{Finger Tapping / 点弦}
\label{ch:tapping}

\eng{Finger tapping is considered a staple of math rock guitar techniques by many. It is widely
employed in the style, and in similar styles that are often attributed to math rock (Midwest emo,
math core, neo soul, Nu Prog and so on). The technique is employed somewhat differently
when compared to more conventional forms of finger tapping, and this difference is what
makes it unique to mathy styles of guitar playing.}

\chn{点弦被许多人认为是数学摇滚吉他技巧的标志。它在这一风格以及与数学摇滚相关的类似风格中被广泛使用（例如 Midwest emo、math core、neo soul、Nu Prog 等）。与更传统的点弦方式相比，数学摇滚点弦的使用方式有所不同，而这种差异恰恰使它在“数学”吉他演奏中独树一帜。}

\eng{I can only speculate on its rise in prominence, perhaps it is due to wanting to play rhythm and
lead guitar at the same time, or it could simply just be to look flashy. Either way, it can sound
excellent when done right and look hella (pun intended) impressive too! It has developed into
being as synonymous with the style as much as the use of odd time signatures has come to be.}

\chn{我只能猜测它走红的原因。也许是吉他手希望同时演奏节奏和主音，或许仅仅是为了看起来更炫。不管怎样，只要演奏得当，它听起来就能非常出色，而且看起来也很惊艳！它已经如同奇数拍号一样，成为这种风格的代名词之一。}

\eng{Here is a list of a few guitarists whom of which regularly employ the technique in their writing
and playing:}

\chn{下面列举了一些在创作和演奏中经常使用点弦技术的吉他手：}

\begin{itemize}[leftmargin=*]
  \item \eng{Yvette Young (Covet)}\chn{Yvette Young（Covet）}
  \item \eng{Tim Collis (TTNG)}\chn{Tim Collis（TTNG）}
  \item \eng{Brock Benzel (Invalids)}\chn{Brock Benzel（Invalids）}
  \item \eng{Dylan Mattheisen (Tiny Moving Parts)}\chn{Dylan Mattheisen（Tiny Moving Parts）}
  \item \eng{Marcos Mena (Standards)}\chn{Marcos Mena（Standards）}
  \item \eng{Tim Henson (Polyphia)}\chn{Tim Henson（Polyphia）}
  \item \eng{Nick Reinhart (Tera Melos)}\chn{Nick Reinhart（Tera Melos）}
\end{itemize}

\eng{In the following section, you will find exercises, technique pointers, and handy tips to get you
on your way to tapping like a mathematical pro!}

\chn{在接下来的部分，你会找到练习、技术提示和实用建议，帮助你像数学摇滚高手一样掌握点弦！}

\paragraph{Finger Tapping Technique Types Explored}
\eng{For this segment we will be exploring two handed tapping and tapping with chords.}
\chn{本节我们将探索双手点弦与和弦点弦两种技巧。}

\eng{Tablature key so you know which hand to be using for tapping in each exercise:}
\chn{下面的吉他谱符号将告诉你在每个练习中哪只手负责点弦：}

\begin{itemize}[leftmargin=*]
  \item \eng{A symbol indicates tapping with your fretting hand (hammer-on from nowhere).}
  \chn{某种符号表示由按弦手点弦（从空中敲击产生声音）。}
  \item \eng{A different symbol indicates tapping with your strumming hand.}
  \chn{另一种符号表示由弹弦手点弦。}
\end{itemize}

\paragraph{Finger Tapping Technique Pointers}
\eng{Evenness of tapping: try to keep volumes even as possible at first. Later, you can experiment
with dynamics in your tapping, but first I recommend trying to be as dynamically even as
possible because this will improve your control.}

\chn{点弦音量均衡：刚开始时尽量保持音量一致。之后你可以尝试在点弦中加入力度变化，但我建议先以动态均匀为目标，因为这会提升你的控制力。}

\paragraph{Developing finger strength balance}
\eng{To achieve better finger strength balance, I recommend dedicating more time to your ring and
little fingers on both hands. The exercise below is an excellent starting point for improving
finger evenness and building strength.}

\chn{为了获得更好的手指力量平衡，我建议双手都多练习无名指和小指。下面的练习是提高手指均衡性和力量的绝佳起点。}

\begin{enumerate}[leftmargin=*]
  \item \eng{Position your fingers on the fifth to eighth frets of the fretboard.}
  \chn{将手指放在指板第 5 到第 8 品。}
  \item \eng{Use your fretting hand to tap the strings in this order: index, middle, ring, and little finger.}
  \chn{用按弦手按顺序点弦：食指、中指、无名指、小指。}
  \item \eng{Move to the ninth to twelfth frets and use your strumming hand to tap in the same finger
    order.}
  \chn{移到第 9 到第 12 品，用弹弦手按同样的指序点弦。}
  \item \eng{A useful tip for your little finger on your strumming hand is to place your ring finger on top
    of it (apply pressure to your little finger's nail) to add more strength to the tap.}
  \chn{对于弹弦手的小指，一个实用技巧是让无名指覆盖在它上面（对小指指甲施加压力），以增强点弦力度。}
  \item \eng{Once you feel comfortable on one string, try moving up and down all the strings.}
  \chn{当你在一根弦上已经熟练后，尝试在所有弦上上下移动练习。}
  \item \eng{Ensure you tap each note for the specified duration (eighth notes: count 1 \& 2 \& 3 \& 4 \&), and
    make sure each note sounds correctly.}
  \chn{确保你每个音都按照指定时值点出（八分音符：数 1 \& 2 \& 3 \& 4 \&），并保证每个音都能清晰发声。}
  \item \eng{Do not forget to use a metronome to help with timing and consistency.}
  \chn{别忘了使用节拍器，它能帮助你保持节奏和稳定。}
\end{enumerate}

\paragraph{Finger Strength \& Dexterity Development}
\eng{As you practice finger tapping, your tapping dexterity will naturally improve. To get you started
though, I suggest the following exercises. These also double up as excellent warm-up routines!}

\chn{随着你练习点弦，你的灵活性会自然提高。然而为了帮助你起步，我建议先做以下练习，它们也非常适合作为热身。}

\eng{The following exercise is adapted from a John Petrucci exercise a guitar teacher showed me
years ago. It has some good stretches and notice how the exercise becomes inverse of itself as
it progresses and arrives back at the original fingering at the beginning of the repeat. I have
added which fingers you should be using next to each note head in the exercise. Here is a key
for which numbers refer to which finger on your fretting hand:}

\chn{下面的练习改编自多年前一位吉他老师给我的 John Petrucci 练习。它包含了很好的伸展动作，并且随着进行会逐渐倒转，最终回到重复开始时的指型。我在谱例旁边标注了每个音符应使用的手指。以下是按弦手手指对应的数字说明：}

\eng{Tip: when tapping, use the underside of your index finger to mute unwanted string ring (when
you are not using it to tap of course!) by lightly touching strings that could potentially ring. To
clarify, you are not to push down on the strings, just apply a light pressure for it to be enough
to stop unwanted strings from ringing out. This is something you should develop if you have
not already because it will make your tapping sound much cleaner. Also, when you are tapping
on higher strings, G, B, or E for example, you should lightly rest the palm of your strumming
hand on the lower strings (E, A, D) to mute them.}

\chn{提示：点弦时，用食指的下侧轻触不需要发声的弦来阻止它们共鸣（当然不是按下，只需轻触即可）。这会让你的点弦听起来更干净。如果你在高音弦上点弦，例如 G、B、E，可以让弹弦手的手掌轻轻搁在低音弦（E、A、D）上以制音。}

\eng{For our strumming hand, we have the same exercises but higher up the fretboard as this is
more common territory for that hand. It will feel awkward at first, and you will probably hit
some bum notes, nevertheless, persevere, it will get better, trust me!}

\chn{对于弹弦手，我们也有一套类似的练习，只不过位置更靠近指板的高把位，因为这是弹弦手更常见的活动区域。一开始会感觉不自然，也可能会击到一些哑音，但坚持下去，它会变得更好，相信我！}

\paragraph{Hand/Wrist Positioning for Finger Tapping}
\eng{To mitigate any unnecessary strain on your fingers, hands, and wrists I have a few tips to help!}

\chn{为了减轻手指、手和手腕的不必要压力，我有几个建议可以帮助你！}

\eng{First, try to keep your fretting hand's wrist straight as possible, it is not always possible but aim
to keep it so:}

\chn{首先，尽量保持按弦手的手腕尽可能笔直。这并不总是可行，但应作为目标。}

\eng{Lastly, for your strumming hand, try to get your hand in-line with your fretboard (see left
picture below for clarification) so that your hand and fingers are facing towards the nut of your
guitar. You will find this more comfortable and easier to play than having you hand and fingers
perpendicular to the fretboard. I like to anchor my thumb on the top side of the fretboard when
finger tapping for stability, but this is not always an option (see the picture on the right below
for clarification).}

\chn{最后，对于弹弦手，尽量让手与指板保持同一直线（手和手指朝向琴枕）。你会发现这比让手和指板垂直更舒适，也更容易弹奏。我喜欢在点弦时将大拇指固定在指板上方以增加稳定性，但这并非总是可行。}

\paragraph{Two Handed Tapping \& Alternating Finger Tapping}
\eng{Two handed tapping mainly consist of alternating tapping (alternating between each hand
when tapping). This is the most prominent version of tapping employed by math rock guitarists.}

\chn{双手点弦主要由交替点弦构成（两只手交替点弦）。这是数学摇滚吉他手最常使用的点弦形式。}

\eng{Occasionally, you might encounter a situation where you need to tap notes simultaneously with
both hands, which, in my opinion, can be quite challenging! For instance, when attempting this
section in my finger-tapping cover of "Santa Claus is Coming to Town," I personally find it more
difficult to coordinate tapping with both hands simultaneously. A helpful tip to make this
exercise easier is to first learn each part separately and then gradually combine the two.}

\chn{偶尔你可能会遇到需要两只手同时点弦的情况，这在我看来相当具有挑战性。例如，在我翻弹《Santa Claus is Coming to Town》的点弦部分时，我个人觉得同时协调两只手更难。一个实用技巧是先分别学习每个声部，然后再逐渐把它们合在一起。}

\eng{Below, you will find a bunch of exercises aimed at improving your two-handed tapping
technique. In later sections, we will explore chords, arpeggios, and scales to create finger
tapping ideas, but for now let's concentrate on getting better at the technique.}

\chn{下面你会找到一系列旨在提升双手点弦技巧的练习。后续章节我们会探索和弦、琶音和音阶以创造点弦素材，但现在我们先集中练好基础技巧。}

\eng{Exercise 1.E goes easy on the strumming hand and has a small ascending and descending
pattern in it.}

\chn{练习 1.E 对弹弦手比较友好，包含一个小的上行和下行模式。}

\eng{Exercise 1.F is built around a major seven arpeggio. It starts with two notes and adds an extra
note in each bar that follows. This is a good exercise to demonstrate simply how you can break
an arpeggio into parts for tapping ideas.}

\chn{练习 1.F 以大七琶音为基础。从两个音开始，每个小节加入一个额外音。这个练习很好地演示了如何将琶音拆分成点弦素材。}

\eng{Exercise 1.G is based around extended seven chord arpeggios. I have split each arpeggio
between both hands. As each bar progresses a new arpeggio of the major scale is introduced.
Take your time with each bar, then glue it all together when you are comfortable with each
individual arpeggio.}

\chn{练习 1.G 基于扩展七和弦琶音。我将每个琶音分配给双手。随着每个小节的进行，会引入一个新的大调琶音。先慢慢练习每个小节，熟练后再把它们连起来。}

\eng{I have included the following two exercises (NS1 \& NS2) as next step exercises. They are more
like the type of tapping you are to come across, but they still have an exercise feel to them.
Never-the-less, I recommend them as the next step to take, and they will set you up nicely to
attempt the test exercises at the end of this section.}

\chn{我将以下两个练习（NS1 和 NS2）作为下一步练习。它们更接近你实际会遇到的点弦形式，但仍保留练习的感觉。无论如何，我建议你把它们作为下一阶段，并以此为基础去挑战本节末尾的测试练习。}

\eng{Test Yourself: Two Handed Finger Tapping Exercises in Context}

\chn{自测：双手点弦在音乐语境中的练习}

\eng{As you have reached the end of the two-hand finger tapping exercises, take a moment to
reflect on your progress so far. Challenge yourself by learning the following two exercises. By
mastering these, you are already ahead of many other guitarists who either have not learned or
have put minimal effort into finger tapping. Pause for a moment, and give yourself a well-
deserved pat on the back, you have earned it.}

\chn{当你完成双手点弦练习后，花一点时间回顾一下迄今为止的进步。挑战自己，学习下面这两个练习。只要掌握它们，你就已经领先于许多尚未学习或付出很少努力的吉他手了。停下来，给自己一个应得的肯定，你配得上这份成就。}

\eng{To round off the two-hand tapping segment, I have listed a bunch of songs that use two-
handed tapping at some point within the song. They are a great way to expand your tapping
repertoire, and through osmosis, you will start to include things that you learn into your own
playing and song writing through already established math rock content.}

\chn{为了结束双手点弦部分，我列出了一些在歌曲中某个时刻使用双手点弦的作品。它们是扩展点弦曲目库的好途径，通过潜移默化，你将把学到的东西融入自己的演奏和创作之中。}

\begin{itemize}[leftmargin=*]
  \item \eng{Maps \& Atlases — Everyplace is a house (Beginner)}\chn{Maps \& Atlases — Everyplace is a house（入门）}
  \item \eng{Tiny Moving Parts — Always Focused Intro (Intermediate)}\chn{Tiny Moving Parts — Always Focused Intro（中级）}
  \item \eng{Don Cabellero — Fire Back About Your New Baby's Sex (Intermediate)}\chn{Don Cabellero — Fire Back About Your New Baby's Sex（中级）}
  \item \eng{Tera Melos — Last Smile for Jaron Intro (Intermediate)}\chn{Tera Melos — Last Smile for Jaron Intro（中级）}
  \item \eng{Tera Melos — 40 Rods to The Hogs Head Intro (Advanced)}\chn{Tera Melos — 40 Rods to The Hogs Head Intro（高级）}
  \item \eng{Covet — Pelagic Intro (Advanced)}\chn{Covet — Pelagic Intro（高级）}
  \item \eng{Standards — Kale \& Strawberries Intro (Advanced)}\chn{Standards — Kale \& Strawberries Intro（高级）}
\end{itemize}

\paragraph{Finger Tapping with Chords}
\eng{This technique is basically what the name implies; strum or arpeggiate a chord and tap notes
along with it. You could also fret a chord or a particular set of notes, but not pick or strum the
shape, and instead tap notes and pull off letting the fretted note or notes of your chord ring
out underneath.}

\chn{这种技巧基本上就是其名称所指：拨响或琶音和弦，同时加入点弦。你也可以按住一个和弦或特定音型，不拨弦，而是点弦并拉开，让按住的音符在下方继续共鸣。}

\eng{This is my favourite tapping technique. It is particularly useful for bands with one guitarist
because it fills out the sound by potentially playing harmony and a melody at the same time.
The technique mainly consists of tapping single notes, or combinations of two, and sometimes,
though rarely, three notes at the same time then pulling off, which lets the notes of the chord
you are fretting ring out as if they were open notes. There are no hard rules on which fingers to
use for tapping, but generally you will feel most comfortable using your index and middle
fingers when not using a pick, and if you are using a pick, try your middle and ring fingers.}

\chn{这是我最喜欢的点弦技巧。对于只有一个吉他手的乐队特别有用，因为它可以同时演奏和声与旋律，丰富声音。该技巧主要包括点单音、两个音的组合，有时（虽少见）还会同时点三个音，然后拉开，让按住的和弦音像开放音一样持续响起。点弦并没有硬性规定使用哪些手指，但通常不使用拨片时更习惯用食指和中指，使用拨片时可以尝试用中指和无名指。}

\eng{Here is an exercise to get you started. Hold an A major barre chord, strum the chord and then
use two fingers to tap and pull off letting the notes of the A major chord ring out underneath.}

\chn{下面是一个入门练习。按住 A 大调把位，拨响和弦，然后用两个手指点弦并拉开，让 A 大调和弦的音在下方持续共鸣。}

\eng{For exercise 1.L, I adapted an exercise from the amazingly talented guitarist Ichika Nito. For this
exercise, I recommend practicing the chord changes first (I have listed them above the exercise
for clarity), and then attempt the tapping parts. This way you are not trying to master two
things at once, which will make the exercise more manageable. Remember to anchor your
thumb on the side of the fret board. Then, when you are comfortable, try doing the same
exercise with a pick and tap the tapped parts with your middle and ring finger.}

\chn{对于练习 1.L，我改编自天才吉他手 Ichika Nito 的练习。建议先练习和弦变化（我在练习上方列出了清晰的和弦），然后再尝试点弦部分。这样你不会同时掌握两件事，练起来更容易。记得将大拇指固定在指板侧面。熟练后，可以尝试使用拨片，并用中指和无名指点弦。}

\eng{By practicing both tapping approaches, you will develop your tapping technique for both
hybrid picking and fingerpicking techniques. Lastly, notice that it is in 5/4 time, so there are five
beats instead of the usual four, and notice you are tapping out a triplet rhythm on each beat
after strumming each chord.}

\chn{通过练习这两种点弦方式，你将同时提升混合拨弦和指弹点弦技巧。最后请注意，练习采用 5/4 拍，因此每小节有五拍而不是通常的四拍，并且每次拨完和弦后，你会用三连音节奏来点弦。}

\eng{For exercise 1.M, fret and hold the chord shapes then tap and pull-off. The tapped notes are
octaves of each chord. This octave higher tapping method is one option you can employ for
finger tapping with chords, but it in context it will be better to mix it up with other notes to
make it sound less like an exercise. Again, I have included the chord shapes for clarity:}

\chn{对于练习 1.M，按住和弦形状然后点弦并拉开。点弦音是每个和弦的同名八度。这种高八度点弦方法是和弦点弦的一种方式，但在乐句中最好与其他音混合使用，避免听起来太像练习。再次说明，我已为你提供了和弦形状以便理解。}

\eng{Test Yourself: Finger Tapping with Chords Exercises in Context}

\chn{自测：和弦点弦在语境中的练习}

\eng{Lastly, let us conclude this section with two examples of how finger tapping with chords can be
used in context.}

\chn{最后，让我们用两个示例来结束本节，展示和弦点弦如何在语境中使用。}
